\documentclass[11pt]{article}

% ---------------------------------------------------------------------------
% Panduan Kode: Quorum-Aware Disruption-Budget Controller (M7)
% Dibangun dengan pdflatex (BUKAN xelatex). Tidak memakai fontspec/babel supaya
% pasti ter-build di mesin mana pun. Seluruh ISI berbahasa Indonesia, untuk
% mahasiswa S1 tahun ke-3 Teknik Informatika.
% ---------------------------------------------------------------------------

\usepackage[T1]{fontenc}
\usepackage[utf8]{inputenc}
\usepackage{amsmath}
\usepackage[margin=1in]{geometry}
\usepackage{parskip}
\usepackage{xcolor}
\usepackage{graphicx}
\usepackage{booktabs}
\usepackage{enumitem}
\usepackage{longtable}
\usepackage{array}
\usepackage{listings}
\usepackage[colorlinks=true,linkcolor=NavyBlue,urlcolor=NavyBlue,citecolor=NavyBlue]{hyperref}

% Warna -----------------------------------------------------------------------
\definecolor{NavyBlue}{RGB}{31,59,115}
\definecolor{CodeBg}{RGB}{246,247,249}
\definecolor{CodeKw}{RGB}{20,80,160}
\definecolor{CodeStr}{RGB}{40,120,40}
\definecolor{CodeCom}{RGB}{120,120,120}
\definecolor{RuleGray}{RGB}{200,200,200}

% Gaya listing umum -----------------------------------------------------------
\lstset{
  basicstyle=\ttfamily\footnotesize,
  backgroundcolor=\color{CodeBg},
  frame=single, rulecolor=\color{RuleGray}, framesep=4pt,
  breaklines=true, breakatwhitespace=false,
  postbreak=\mbox{\textcolor{RuleGray}{$\hookrightarrow$}\space},
  columns=fullflexible, keepspaces=true, showstringspaces=false,
  tabsize=2, captionpos=b, xleftmargin=10pt, xrightmargin=4pt,
  aboveskip=8pt, belowskip=8pt,
}

% Bahasa Elixir untuk listings ------------------------------------------------
\lstdefinelanguage{elixir}{
  morekeywords={def,defp,defmodule,defmacro,defstruct,do,end,fn,case,cond,
    if,else,unless,when,and,or,not,in,for,with,receive,after,try,catch,raise,
    use,alias,import,require,true,false,nil,@impl,@moduledoc,@doc,@spec,@type},
  otherkeywords={|>,->,<-,=>,::,\&},
  sensitive=true, morecomment=[l]{\#}, morestring=[b]", morestring=[b]',
  keywordstyle=\color{CodeKw}\bfseries, commentstyle=\color{CodeCom}\itshape,
  stringstyle=\color{CodeStr}, identifierstyle=\color{black},
}
\lstdefinelanguage{yaml}{
  keywords={true,false,null,name,kind,apiVersion,metadata,spec,template,containers,
    image,ports,env,value,replicas,selector,strategy,rollingUpdate,maxUnavailable,
    maxSurge,minAvailable,namespace,labels,rules,apiGroups,resources,verbs,
    roleRef,subjects,serviceAccountName,clusterIP,readinessProbe,livenessProbe,httpGet,path,port},
  sensitive=true, comment=[l]{\#}, morestring=[b]", morestring=[b]',
  keywordstyle=\color{CodeKw}\bfseries, commentstyle=\color{CodeCom}\itshape,
  stringstyle=\color{CodeStr},
}
\lstdefinestyle{ex}{language=elixir}
\lstdefinestyle{yml}{language=yaml}
\lstdefinestyle{sh}{language=bash,keywordstyle=\color{CodeKw}\bfseries,
  commentstyle=\color{CodeCom}\itshape,stringstyle=\color{CodeStr}}

\newcommand{\code}[1]{\texttt{\small #1}}
\newcommand{\istilah}[1]{\textbf{#1}}
\setlist{topsep=3pt,itemsep=2pt,parsep=0pt}

% =============================================================================
\title{\textbf{Panduan Kode (untuk Mahasiswa)}\\[4pt]
\large Quorum-Aware Disruption-Budget Controller:\\
Menjaga Kuorum Cluster Elixir/BEAM saat Pemeliharaan Kubernetes}
\author{Dokumentasi teknis proyek \code{02\_quorum\_aware\_disruption\_budget}}
\date{Dihasilkan otomatis --- versi panduan kode}

\begin{document}
\maketitle

\begin{abstract}
\noindent
Dokumen ini menjelaskan \textbf{kode} sebuah \emph{controller lintas-lapis}: sistem yang membuat
Kubernetes ``tahu'' berapa banyak pod boleh dimatikan sekaligus tanpa memecah \emph{kuorum} sebuah
cluster Elixir/BEAM. Tujuannya menutup apa yang kami sebut \emph{quorum gap}: celah antara batas
gangguan yang disetel manual di Kubernetes (\emph{PodDisruptionBudget}) dan jumlah anggota minimum yang
sebenarnya dibutuhkan aplikasi agar tetap berwenang. Panduan ini ditulis dengan \textbf{bahasa yang
mudah} untuk mahasiswa S1 tahun ke-3: tiap istilah teknis dijelaskan saat pertama muncul, disertai
analogi sehari-hari, cuplikan kode asli, dan \textbf{langkah persis} menjalankan seluruh eksperimen
(RQ1--RQ10). Semua nama fungsi, file, dan angka diambil langsung dari kode di
\code{02\_quorum\_aware\_disruption\_budget/}.
\end{abstract}

\tableofcontents
\bigskip \hrule \bigskip

% =============================================================================
\section{Pendahuluan: Masalah ``Quorum Gap''}
\label{sec:intro}

\subsection{Gambaran singkat}
Bayangkan sebuah aplikasi yang berjalan sebagai \emph{cluster}: banyak salinan (pod) yang saling
terhubung dan harus \emph{sepakat} untuk bekerja benar. Banyak sistem terdistribusi membutuhkan
\istilah{kuorum} (quorum): jumlah anggota hidup minimum --- biasanya \emph{mayoritas} --- agar cluster
tetap berwenang mengambil keputusan. Kalau anggota hidup turun di bawah kuorum, cluster bisa
\emph{pecah} (dua bagian sama-sama merasa berwenang $=$ \istilah{split-brain}) atau kehilangan salinan
state yang sah.

Masalahnya: ketika Kubernetes melakukan \emph{pemeliharaan} (rolling update, drain node, scale-down),
ia mematikan beberapa pod sekaligus. Berapa banyak yang boleh mati bersamaan diatur oleh sebuah aturan
bernama \istilah{PodDisruptionBudget (PDB)}. Tetapi PDB \emph{disetel manual} dalam satuan jumlah pod
dan \emph{tidak tahu} berapa kuorum aplikasi. Jadi pemeliharaan yang ``terlihat aman'' di tingkat pod
bisa diam-diam mematikan terlalu banyak anggota dan memecah kuorum --- padahal semua probe pod masih
hijau. Itulah \emph{quorum gap}. Controller di proyek ini menutupnya dengan \textbf{menurunkan PDB dari
kuorum runtime}, sehingga pemeliharaan tak pernah memecah kuorum.

\subsection{Apa itu gangguan sukarela \& PDB?}
Kubernetes membedakan dua jenis gangguan. \istilah{Gangguan tak-sukarela} (involuntary): mesin mati,
kernel panik, OOM --- tak ada peringatan. \istilah{Gangguan sukarela} (voluntary): rolling update,
\code{kubectl drain}, scale-down --- direncanakan operator. Untuk gangguan sukarela, Kubernetes
menyediakan satu penjaga: PDB. PDB menyatakan \code{minAvailable} (paling sedikit pod yang harus tetap
\emph{Ready}) atau \code{maxUnavailable} (paling banyak yang boleh mati). API \emph{eviction} dan alur
\code{kubectl drain} menghormati PDB: eviction yang melanggarnya \emph{ditolak}.

\subsection{Apa itu kuorum di BEAM?}
Di Elixir/BEAM, pod-pod membentuk cluster lewat \code{libcluster}. State dibagi lewat berbagai cara
yang kebutuhannya berbeda: \code{:global} (nama unik se-cluster; kalau partisi pulih, salah satu
duplikat dibunuh --- state-nya hilang), \code{Horde} (registry/supervisor terdistribusi), atau library
konsensus (butuh \emph{mayoritas} ketat untuk maju). Yang sama dari semuanya: ada ambang \emph{kuorum}
$Q$ di bawah mana cluster berhenti berwenang atau berisiko split-brain. $Q$ adalah properti yang
\emph{aplikasi} tahu (mayoritas, atau \emph{floor} yang dideklarasikan operator), \emph{Kubernetes
tidak}.

\subsection{Mengapa PDB tetap (statis) itu salah}
\label{sec:why-static}
Nilai aman PDB \textbf{bukan konstanta}. Di bawah kuorum mayoritas, $Q$ mengikuti ukuran cluster yang
\emph{diinginkan}: mayoritas dari 5 adalah 3, dari 9 adalah 5. Saat operator menaikkan/menurunkan
jumlah replika, $Q$ berubah. Sebuah \code{minAvailable} statis karena itu bisa terlalu kecil untuk
sebagian ukuran (mengizinkan gangguan pemecah-kuorum) atau terlalu besar untuk sebagian lain
(menghambat pemeliharaan yang sebenarnya aman). Tak ada satu angka yang benar untuk semua ukuran.

\istilah{Kunci desain (penting!):} kuorum harus dipatok ke ukuran cluster yang \textbf{diinginkan}
(jumlah replika Deployment), \textbf{bukan} ukuran \emph{live} sesaat. Kalau memakai ukuran live, sebuah
partisi yang menyusut akan selalu mengira ``mayoritas dari ukuranku sendiri'' terpenuhi --- yakni bug
split-brain yang justru ingin dicegah. Satu baris di \code{Cluster.quorum\_threshold/1} (Bagian~\ref{sec:mod-cluster})
menegakkan ini, dan eksperimen RQ10 mengukurnya.

\subsection{Apa yang dilakukan controller ini}
Tiap pod menjalankan \emph{probe} read-only yang melaporkan ukuran cluster + ambang kuorum $Q$. Sebuah
\emph{operator} (pod terpisah) membaca probe semua pod, menghitung budget, lalu menambal
\code{minAvailable} PDB ke $Q$. Setelah itu, eviction yang akan menurunkan anggota di bawah $Q$ ditolak
oleh Kubernetes. Rumus budget-nya:
\begin{equation}
\code{minAvailable} = \max(Q,\,Q_{\min}), \qquad
\code{maxUnavailable} = \min(N - \code{minAvailable},\ c),
\label{eq:budget}
\end{equation}
dengan $N$ ukuran (diinginkan), $Q_{\min}$ batas bawah, dan $c$ kapasitas handoff (pembatas opsional).

\subsection{Analogi sehari-hari}
Bayangkan rapat dewan beranggota 5 orang yang aturannya: keputusan sah hanya bila \emph{mayoritas}
(3 orang) hadir. Itu ``kuorum''-nya, $Q=3$. Suatu hari panitia gedung ingin merenovasi ruangan dan
berkata ``boleh 3 orang izin sekaligus, kan masih ada sisanya?'' --- terdengar wajar. Tapi kalau 3
orang benar-benar izin bersamaan, yang tersisa hanya 2 $<$ 3: rapat \emph{tak lagi sah}, padahal tiap
orang yang izin merasa ``cuma satu orang, tak masalah''. Panitia gedung (Kubernetes) tak tahu aturan
kuorum dewan (aplikasi). Controller ini ibarat sekretaris yang memberi tahu panitia: ``maksimal 2 orang
boleh izin sekaligus'' --- diturunkan dari aturan kuorum yang sebenarnya.

\subsection{Contoh konkret (dengan angka)}
\label{sec:contoh}
Cluster 5 pod, kuorum mayoritas $Q = \lfloor 5/2\rfloor + 1 = 3$.
\begin{itemize}[leftmargin=1.4em]
  \item \textbf{PDB statis (tak sadar kuorum), \code{maxUnavailable}=3:} rolling update mengusir 3 pod
        sekaligus $\rightarrow$ tersisa $5-3 = 2 < 3$. \textbf{Kuorum PECAH.} (Inilah yang RQ1 ukur:
        anggota tersedia jatuh ke 2.)
  \item \textbf{Quorum-aware, \code{minAvailable}=$Q$=3 (jadi \code{maxUnavailable}$=5-3=2$):} paling
        banyak 2 pod turun sekaligus $\rightarrow$ tersisa $5-2 = 3 = Q$. \textbf{Kuorum AMAN.}
\end{itemize}
Bedanya cuma satu angka (2 vs 3), tapi itu garis antara cluster yang utuh dan cluster yang pecah.

% =============================================================================
\section{Struktur Folder Proyek}
\label{sec:struktur}
\begin{lstlisting}[style=sh]
02_quorum_aware_disruption_budget/
|-- paper/                 # makalah LaTeX (Wiley) -- bukan bagian kode
|-- code/
|   |-- app/               # aplikasi Elixir/OTP (inti)
|   |   |-- lib/quorum_budget/
|   |   |   |-- quorum.ex          # POLICY murni: hitung budget (Persamaan 1)
|   |   |   |-- cluster.ex         # baca keanggotaan + turunkan ambang kuorum Q
|   |   |   |-- quorum_probe.ex    # sensor read-only (GenServer)
|   |   |   |-- probe_http.ex      # endpoint HTTP /probe (JSON)
|   |   |   |-- pdb_operator.ex    # operator: tambal minAvailable PDB via kubectl
|   |   |   |-- quorum_workload.ex # beban kerja bergantung-kuorum (RQ8)
|   |   |   |-- disruptor.ex       # mesin eksperimen: rolling update + ukur
|   |   |   |-- harness.ex         # driver pengukuran tiap RQ
|   |   |   `-- application.ex     # pohon supervisi OTP (pilih peran app/operator)
|   |   |-- harness/        # skrip eksperimen (run/policy/workload/reactive .exs)
|   |   |-- test/           # tes unit (quorum_test) + cluster (cluster_test)
|   |   `-- config/         # config.exs, test.exs, runtime.exs, ...
|   |-- k8s/               # manifes Kubernetes + rq7_eviction.sh
|   |-- analysis/plot.py   # render figur dari data/*.csv
|   `-- Dockerfile
|-- data/                  # results_*.csv (HASIL PENGUKURAN NYATA)
|-- figures/               # eval_*.pdf (dirender plot.py)
`-- code_guide/            # dokumen ini
\end{lstlisting}
Tiga folder yang paling sering dilihat: \code{lib/quorum\_budget/} (kode inti), \code{harness/} (cara
menjalankan eksperimen), dan \code{data/} (hasil mentah yang dipakai paper).

% =============================================================================
\section{Library yang Dipakai dan Perannya}
\label{sec:library}
Sengaja \emph{minimal} --- nilai mekanisme ada pada \emph{dari mana} budget berasal, bukan pada mesin
yang berat.
\begin{longtable}{@{}p{0.20\linewidth} p{0.14\linewidth} p{0.58\linewidth}@{}}
\toprule
\textbf{Library} & \textbf{Versi} & \textbf{Peran} \\
\midrule \endhead
\code{libcluster} & 3.5 & Membentuk cluster BEAM di Kubernetes (pod menemukan peer lewat API server). \\
\code{bandit}     & 1.12 & Server HTTP kecil untuk menyajikan endpoint \code{/probe}. \\
\code{plug}       & 1.20 & Router HTTP (\code{Plug.Router}) di atas Bandit. \\
\code{jason}      & 1.4 & Encode/decode JSON (berbicara dengan API Kubernetes lewat \code{kubectl}). \\
\bottomrule
\end{longtable}
Catatan: M7 berfokus pada keanggotaan-cluster + kuorum, jadi \textbf{tidak} memakai Horde (registry
handoff) --- berbeda dari proyek 01. Hal-hal terdistribusi memakai modul bawaan Erlang
(\code{Node}, \code{:rpc}).

% =============================================================================
\section{Penjelasan Modul-per-Modul}
\label{sec:modul}
Gambar~\ref{fig:seqmod} memperlihatkan bagaimana modul-modul bekerja sama saat satu rolling update
(eksperimen) berjalan.

\begin{figure}[h]\centering
  \includegraphics[width=0.95\linewidth]{fig/seq_modules.pdf}
  \caption{Interaksi antar-modul saat satu rolling update diukur (Disruptor menggerakkan
  eviction/measure/rejoin; Cluster \& QuorumWorkload memberi pembacaan nyata).}\label{fig:seqmod}
\end{figure}

\subsection{\texttt{QuorumBudget.Quorum} --- kebijakan (policy) murni}
\label{sec:mod-quorum}
\textbf{Tugas:} menghitung budget PDB dari ukuran cluster $N$ dan ambang kuorum $Q$. \textbf{Sifat:}
murni (tanpa proses, tanpa I/O) --- inilah bagian yang dibuktikan paper dan diuji langsung.
\begin{lstlisting}[style=ex]
def budget(reading, opts \\ []) do
  n = reading.n;  q = reading.q
  q_min = Keyword.get(opts, :q_min, 1)
  cap = Map.get(reading, :cap, n)                 # batas opsional dari kapasitas handoff
  min_available  = q |> max(q_min) |> min(n)      # = kuorum: menegakkan A >= Q
  max_unavailable = (n - min_available) |> min(trunc(cap)) |> max(0)
  %{min_available: min_available, max_unavailable: max_unavailable}
end
\end{lstlisting}
Tiga fungsi pendukung melengkapinya. Perhatikan betapa pendeknya --- ini ``otak'' mekanisme, dan
sengaja sederhana supaya bisa dibuktikan benar:
\begin{lstlisting}[style=ex]
# Kuorum mayoritas untuk cluster ukuran n (3 dari 5, 4 dari 7, ...)
def majority(n) when is_integer(n) and n > 0, do: div(n, 2) + 1

# Apakah mengusir d pod dari a anggota tersedia akan melanggar kuorum q?
def violates_quorum?(a, d, q), do: a - d < q

# Keputusan untuk SATU eviction: izinkan hanya bila setelah satu pod pergi,
# anggota tersedia masih >= min_available. Inilah yang ditegakkan PDB Kubernetes.
def admit_eviction(available, min_available) do
  if available - 1 >= min_available, do: :allow, else: :deny
end
\end{lstlisting}
Karena \code{budget/2} dan kawan-kawannya murni (input $\rightarrow$ output, tanpa efek samping),
mereka diuji langsung di \code{quorum\_test.exs} tanpa perlu cluster sama sekali.

\subsection{\texttt{QuorumBudget.Cluster} --- baca keanggotaan \& kuorum}
\label{sec:mod-cluster}
\textbf{Tugas:} melaporkan ukuran cluster dan menurunkan $Q$. \code{size/0} $=$
\code{length(peers())+1}. Fungsi terpenting --- dan \emph{kunci keamanan} --- adalah
\code{quorum\_threshold/1}:
\begin{lstlisting}[style=ex]
def quorum_threshold(live \\ size()) do
  desired = Application.get_env(:quorum_budget, :cluster_size, live)  # UKURAN DIINGINKAN
  case Application.get_env(:quorum_budget, :quorum, :majority) do
    :majority -> Quorum.majority(desired)        # mayoritas dari ukuran DIINGINKAN
    q when is_integer(q) and q > 0 -> q          # floor tetap (mis. :global/CRDT)
    _ -> Quorum.majority(desired)
  end
end
\end{lstlisting}
Perhatikan \code{desired}: kuorum dihitung dari ukuran yang \emph{diinginkan}, bukan ukuran \emph{live}.
Inilah satu baris yang mencegah split-brain (lihat Bagian~\ref{sec:why-static}).
\begin{lstlisting}[style=ex]
def size, do: length(peers()) + 1

# Peer-app yang terhubung, MENGECUALIKAN node pengendali (mis. primary@... di harness tes,
# atau pod operator) -- ia bukan anggota cluster aplikasi, jadi tak dihitung sebagai kuorum.
def peers do
  case Application.get_env(:quorum_budget, :control_node_prefix) do
    nil -> Node.list()
    prefix -> Enum.reject(Node.list(), &String.starts_with?(to_string(&1), prefix))
  end
end

# Apakah node INI sedang berkuorum: ukuran yang dilihatnya memenuhi ambang?
def in_quorum? do
  n = size()
  n >= quorum_threshold(n)
end
\end{lstlisting}
\code{in\_quorum?/0} dipakai \code{QuorumWorkload} untuk memutuskan boleh tidaknya \emph{commit}: kalau
node tak melihat kuorum, ia tak boleh maju --- persis seperti grup konsensus nyata.

\subsection{\texttt{QuorumBudget.QuorumProbe} --- sensor read-only}
\textbf{Tugas:} \code{GenServer} (proses server OTP) yang, saat ditanya, mengembalikan map
\code{\%\{n, q, cap, in\_quorum\}}. Tak mengubah apa pun --- hanya membaca. Operator memanggilnya lewat
HTTP untuk mendapat $n$ dan $q$.

\subsection{\texttt{QuorumBudget.ProbeHTTP} --- endpoint HTTP}
\textbf{Tugas:} \code{Plug.Router} kecil yang menyajikan \code{GET /probe} (JSON dari pembacaan probe)
dan \code{GET /healthz} (cek hidup untuk liveness probe Kubernetes).

\subsection{\texttt{QuorumBudget.PDBOperator} --- operator Kubernetes}
\label{sec:mod-operator}
\textbf{Tugas:} koordinator lintas-lapis. Tiap 5 detik: (1) daftar IP pod via \code{kubectl}, (2) baca
\code{/probe} tiap pod, (3) hitung \code{minAvailable} (ambil yang \emph{terbesar} antar-pod --- paling
aman saat pod tak sepakat soal ukuran), lalu (4) tambal \code{spec.minAvailable} PDB. Dibungkus
\code{try/catch}: error API sementara hanya di-log, operator tak crash, budget terakhir tetap dipakai
(fail-safe). Hanya jalan dalam peran \code{operator} (env \code{QUORUM\_ROLE=operator}).
\begin{lstlisting}[style=ex]
def reconcile do
  readings = for ip <- pod_ips(), r = probe(ip), do: r       # probe tiap pod (nil dilewati)
  if readings != [] do
    min_avail =
      readings
      |> Enum.map(fn r -> Quorum.budget(r, q_min: q_min()).min_available end)
      |> Enum.max()                                           # paling aman antar-pod
    patch_pdb(min_avail)                                      # tambal minAvailable PDB
  end
end

defp patch_pdb(min_avail) do
  patch = Jason.encode!(%{spec: %{minAvailable: min_avail}})
  System.cmd("kubectl", ["patch", "pdb", pdb(), "-n", namespace(),
                         "--type", "merge", "-p", patch])
end
\end{lstlisting}

\subsection{\texttt{QuorumBudget.QuorumWorkload} --- beban kerja bergantung-kuorum (RQ8)}
\label{sec:mod-workload}
\textbf{Tugas:} meniru grup konsensus / singleton \code{:global} yang otoritatif. Tiap 25 ms ia cek
\code{Cluster.in\_quorum?}: bila ya, \emph{commit} berhasil; bila tidak, commit \emph{diblokir}. Jadi
ia mengukur \textbf{konsekuensi nyata} kehilangan kuorum: anggota yang masih ``hidup'' menurut probe
pod tetapi tak bisa melayani satu pun permintaan yang butuh kuorum.
\begin{lstlisting}[style=ex]
@tick_ms 25
def handle_info(:tick, s) do
  s =
    if Cluster.in_quorum?() do
      %{s | committed: s.committed + 1}   # berkuorum -> commit berhasil
    else
      %{s | blocked: s.blocked + 1}       # tak berkuorum -> commit DIBLOKIR (progres hilang)
    end
  Process.send_after(self(), :tick, @tick_ms)
  {:noreply, s}
end
\end{lstlisting}
Jadi \code{blocked} pada survivor selama gangguan $=$ ukuran progres yang \emph{hilang} karena kuorum
pecah. Pada quorum-aware, survivor tak pernah keluar kuorum $\rightarrow$ \code{blocked} survivor $=0$.

\subsection{\texttt{QuorumBudget.Disruptor} --- mesin eksperimen}
\label{sec:mod-disruptor}
\textbf{Tugas:} mensimulasikan rolling update \emph{sungguhan} pada cluster peer dan
\emph{mengukur}-nya. Untuk tiap batch (= ukuran budget) ia benar-benar memutus node dari cluster, lalu
mengukur ukuran cluster tersisa dari \code{Node.list/0} node survivor, lalu menyambungkannya kembali.
\begin{lstlisting}[style=ex]
{min_avail, blocked, stalled} =
  app_nodes |> Enum.chunk_every(batch)
  |> Enum.reduce({n, 0, 0}, fn chunk, {min_a, blk, stl} ->
    survivors = app_nodes -- chunk
    Enum.each(chunk, &evict(&1, app_nodes))        # putus batch (Node.disconnect)
    if pace_ms > 0, do: Process.sleep(pace_ms)
    a = measure_available(app_nodes, chunk)        # UKUR ukuran tersisa (nyata)
    stalled_now = count_stalled(survivors)         # survivor yang out-of-quorum
    Enum.each(chunk, &rejoin(&1, app_nodes))        # sambung lagi (pengganti masuk)
    {min(min_a, a), blk + if(a < q, do: 1, else: 0), stl + stalled_now}
  end)
\end{lstlisting}
Metriknya: \code{min\_available} (ukuran terkecil terukur), \code{quorum\_violated?} ($<Q$),
\code{blocked\_survivors} (survivor yang tak bisa commit), \code{maintenance\_ms}.

\subsection{\texttt{QuorumBudget.Harness} --- driver eksperimen}
\label{sec:mod-harness}
\textbf{Tugas:} membungkus pengukuran tiap RQ menjadi baris CSV. Fungsi: \code{spawn\_mesh/1} \&
\code{stop\_mesh/1} (bentuk/hentikan cluster peer), \code{safety/3} (RQ1+RQ2), \code{workload/3} (RQ8),
\code{reactive/2} (RQ9), \code{adaptivity/sensitivity/overhead/scale/flap} (RQ3--RQ6, RQ10, policy
murni).

\subsection{\texttt{QuorumBudget.Application} --- pohon supervisi}
\textbf{Tugas:} titik masuk OTP. Memilih anak supervisi sesuai peran: \code{:operator} $\rightarrow$
\code{PDBOperator}; \code{:app} (default) $\rightarrow$ pembentukan cluster + \code{QuorumProbe} +
\code{QuorumWorkload} + endpoint HTTP. Satu image, dua peran (env \code{QUORUM\_ROLE}).

% =============================================================================
\section{Cara Kerja Langkah-demi-Langkah}
\label{sec:alur}
Ada dua ``alur'' berbeda. \textbf{Alur A} adalah eksperimen terkontrol pada cluster peer lokal (tanpa
Kubernetes) --- inilah yang menghasilkan angka RQ1--RQ6, RQ8--RQ10. \textbf{Alur B} adalah loop
Kubernetes nyata (operator menambal PDB; API menolak eviction) --- RQ7.

\begin{figure}[h]\centering
  \includegraphics[width=0.95\linewidth]{fig/seq_crosslayer.pdf}
  \caption{Alur B: operator membaca probe lalu menambal PDB; saat drain, API menolak eviction yang akan
  memecah kuorum.}\label{fig:seqcross}
\end{figure}

\subsection{Alur A --- eksperimen rolling update lokal}
\begin{enumerate}[leftmargin=1.4em]
  \item \code{Harness.spawn\_mesh(n)} memulai $n$ node peer BEAM dan menghubungkannya full mesh.
        Tiap peer dibawa config \code{cluster\_size = n} (memaku kuorum ke ukuran diinginkan).
  \item Harness memanggil \code{Disruptor.rolling\_cycle} dengan ukuran batch sesuai policy:
        \emph{static} (mis. $\lceil n/2\rceil$, mengabaikan kuorum) atau \emph{quorum-aware}
        ($n-Q$, dari Persamaan~\ref{eq:budget}).
  \item Untuk tiap batch: putus node (eviction) $\rightarrow$ ukur ukuran tersisa + cek survivor
        berkuorum $\rightarrow$ sambung kembali (pengganti).
  \item Harness mengumpulkan metrik dan menulis \code{data/results\_*.csv}.
\end{enumerate}

\subsection{Alur B --- loop Kubernetes nyata}
\begin{enumerate}[leftmargin=1.4em]
  \item Operator (\code{PDBOperator}) tiap 5 detik: \code{kubectl get pods} $\rightarrow$ HTTP
        \code{/probe} tiap pod $\rightarrow$ \code{Quorum.budget} $\rightarrow$
        \code{kubectl patch pdb minAvailable=Q}.
  \item Saat admin menjalankan \code{kubectl drain} / rolling update, API eviction memeriksa PDB:
        bila menurunkan anggota di bawah \code{minAvailable}, eviction \textbf{ditolak} (HTTP 429);
        selainnya diizinkan.
\end{enumerate}

% =============================================================================
\section{Penerapan Kubernetes (Manifes)}
\label{sec:k8s}
Folder \code{code/k8s/} berisi manifes untuk menjalankan sistem di cluster sungguhan. Tiga yang inti:

\textbf{1. PodDisruptionBudget} --- yang ditambal operator. \code{minAvailable} awal di sini hanya
placeholder; operator menimpanya dengan kuorum mayoritas tiap reconcile.
\begin{lstlisting}[style=yml]
apiVersion: policy/v1
kind: PodDisruptionBudget
metadata: {name: quorum, namespace: default}
spec:
  minAvailable: 3            # = majority(5); ditambal operator
  selector: {matchLabels: {app: quorum}}
\end{lstlisting}

\textbf{2. Deployment aplikasi} (\code{app.yaml}, dipotong) --- $N$ replika yang membentuk cluster BEAM,
dengan \code{POD\_IP} dari Downward API dan probe HTTP:
\begin{lstlisting}[style=yml]
spec:
  replicas: 5
  template:
    spec:
      containers:
        - name: quorum
          image: quorum_budget:latest
          env:
            - {name: QUORUM_ROLE, value: "app"}
            - name: POD_IP
              valueFrom: {fieldRef: {fieldPath: status.podIP}}
          readinessProbe: {httpGet: {path: /probe, port: 4000}}
\end{lstlisting}

\textbf{3. RBAC hak-paling-kecil} (\code{rbac.yaml}, dipotong) --- operator HANYA boleh membaca pod dan
membaca/menambal PDB:
\begin{lstlisting}[style=yml]
rules:
  - apiGroups: [""]
    resources: ["pods"]
    verbs: ["get", "list"]
  - apiGroups: ["policy"]
    resources: ["poddisruptionbudgets"]
    verbs: ["get", "list", "patch"]
\end{lstlisting}
Satu image (\code{Dockerfile}: release Elixir + \code{kubectl}) dipakai peran \code{app} dan
\code{operator}, dibedakan env \code{QUORUM\_ROLE}. Memuat ke kind: \code{kind load docker-image
quorum\_budget:latest}.

% =============================================================================
\section{Bagaimana Kode Menjawab RQ1--RQ10}
\label{sec:rq}
\begin{longtable}{@{}p{0.06\linewidth} p{0.30\linewidth} p{0.28\linewidth} p{0.27\linewidth}@{}}
\toprule
\textbf{RQ} & \textbf{Fokus} & \textbf{Kode penghasil} & \textbf{Data / Figur} \\
\midrule \endhead
RQ1 & safety (kuorum pecah?) & \code{run.exs}, \code{cluster\_test.exs} & \code{results\_safety.csv} $\rightarrow$ \code{eval\_safety.pdf} \\
RQ2 & efficiency (waktu maintenance) & \code{run.exs} & \code{results\_efficiency.csv} $\rightarrow$ \code{eval\_efficiency.pdf} \\
RQ3 & adaptivity (budget ikut $N$) & \code{policy.exs} & \code{results\_scale.csv} $\rightarrow$ \code{eval\_adaptivity.pdf} \\
RQ4 & overhead (biaya hitung) & \code{policy.exs} & \code{results\_overhead.csv} \\
RQ5 & robustness (sensitivitas $q$) & \code{policy.exs} & \code{results\_sensitivity.csv} $\rightarrow$ \code{eval\_sensitivity.pdf} \\
RQ6 & scalability ($O(1)$) & \code{policy.exs} & \code{results\_scale.csv} $\rightarrow$ \code{eval\_scale.pdf} \\
RQ7 & cluster nyata (eviction) & \code{k8s/rq7\_eviction.sh} & \code{results\_eviction.csv} \\
RQ8 & realistic workload & \code{workload.exs} & \code{results\_workload.csv} $\rightarrow$ \code{eval\_workload.pdf} \\
RQ9 & vs baseline reactive & \code{reactive.exs} & \code{results\_reactive.csv} $\rightarrow$ \code{eval\_reactive.pdf} \\
RQ10 & robust ke membership flap & \code{policy.exs} & \code{results\_flap.csv} $\rightarrow$ \code{eval\_flap.pdf} \\
\bottomrule
\end{longtable}

\paragraph{RQ1 (safety).} Budget statis vs quorum-aware pada rolling update nyata. Statis
\emph{memecah kuorum di setiap} ukuran (anggota tersedia turun di bawah $Q$); quorum-aware menahan di
$Q$, \textbf{0} pelanggaran.
\paragraph{RQ2 (efficiency).} Budget aman tak harus lambat: quorum-aware (batch $N-Q$) jauh lebih cepat
daripada konservatif (satu-per-satu), keduanya aman.
\paragraph{RQ3 (adaptivity).} \code{minAvailable} yang dihitung melacak mayoritas $\lfloor N/2\rfloor+1$
persis saat $N$ tumbuh.
\paragraph{RQ4 (overhead) \& RQ6 (scalability).} Perhitungan budget hanya beberapa operasi bilangan
bulat: $\sim$50 ns/panggilan, \emph{datar} (O(1)) dari $N=3$ sampai $10\,001$.
\paragraph{RQ5 (robustness).} \code{maxUnavailable} $=N-q$ linear; \emph{under}-estimasi $Q$ memberi
budget tak aman $\Rightarrow$ estimasi harus di-bias ke atas (itulah fail-safe).
\paragraph{RQ7 (cluster nyata).} Di \code{kind}: PDB sadar-kuorum (\code{minAvailable}=$Q$) membuat API
eviction \emph{menolak} eviction pemecah-kuorum; PDB tak-sadar-kuorum mengizinkannya.
\paragraph{RQ8 (realistic workload).} Beban kerja bergantung-kuorum: budget statis memblokir
$\lfloor N/2\rfloor$ survivor (tak bisa commit walau ``Ready''), quorum-aware \textbf{0}.
\paragraph{RQ9 (vs reactive).} Baseline yang menaikkan \code{minAvailable} setelah tiap pelanggaran
memecah kuorum berkali-kali sebelum konvergen; model-based benar dari probe pertama.
\paragraph{RQ10 (membership flap).} Memaku budget ke ukuran \emph{diinginkan} membuatnya stabil + aman
saat ukuran live berkedip; desain naif (patok live) churn PDB \emph{dan} jadi tak-aman.

\subsection{Contoh keluaran nyata}
\label{sec:keluaran}
Angka di bawah dari run sungguhan (di \code{data/*.csv}); semua bisa kamu reproduksi (Bagian~\ref{sec:run}).
\begin{itemize}[leftmargin=1.4em]
  \item \textbf{RQ1} (\code{results\_safety.csv}): \emph{static} memecah kuorum di $N=3,5,7,9$ (anggota
        tersedia minimum $1/2/3/4$ vs $Q$ $2/3/4/5$); \emph{quorum-aware} $0$ pelanggaran.
  \item \textbf{RQ2} (\code{results\_efficiency.csv}): konservatif (satu-per-satu) $\approx$3{,}49 s vs
        quorum-aware $\approx$1{,}21 s $=$ $\sim$2{,}9$\times$ lebih cepat, dua-duanya aman.
  \item \textbf{RQ4} (\code{results\_overhead.csv}): hitung budget $\approx$50 ns/panggilan.
  \item \textbf{RQ8} (\code{results\_workload.csv}): static memblokir $2/3/4$ survivor (di $N=5/7/9$);
        quorum-aware $0$.
  \item \textbf{RQ9} (\code{results\_reactive.csv}): baseline reactive memecah kuorum \textbf{4$\times$}
        ($minAvailable$ $1\to2\to3\to4$) sebelum konvergen di $5$; model-based $0$.
  \item \textbf{RQ10} (\code{results\_flap.csv}): naif $28$ patch PDB $+$ $14$ sampel tak-aman per 100;
        quorum-aware $0/0$.
  \item \textbf{RQ7} (\code{results\_eviction.csv}): di kind, PDB \code{minAvailable}=3 membuat API
        \emph{menolak} (HTTP 429) eviction yang akan menurunkan anggota di bawah 3.
\end{itemize}

% =============================================================================
\section{Panduan Menjalankan Eksperimen}
\label{sec:run}
\textbf{Prasyarat:} Elixir/OTP terpasang; untuk RQ7 perlu cluster \code{kind} hidup. Dari
\code{code/app/}. Catatan: eksperimen cluster memakai \code{:peer} dengan cookie \code{ck};
\textbf{jangan} \code{pkill -f '...@127.0.0.1'} (pola itu cocok dengan shell yang sedang berjalan dan
mematikannya) --- matikan beam yatim berdasarkan PID.

\begin{lstlisting}[style=sh]
# RQ1 (safety) + RQ2 (efficiency) -- cluster peer nyata
MIX_ENV=test elixir --name primary@127.0.0.1 --cookie ck -S mix run harness/run.exs
# -> data/results_safety.csv, results_efficiency.csv

# RQ3-RQ6 + RQ10 (policy murni; tanpa cluster)
MIX_ENV=test mix run harness/policy.exs
# -> results_scale.csv, results_sensitivity.csv, results_overhead.csv, results_flap.csv

# RQ8 (realistic workload)
MIX_ENV=test elixir --name primary@127.0.0.1 --cookie ck -S mix run harness/workload.exs
# -> results_workload.csv

# RQ9 (baseline reactive)
MIX_ENV=test elixir --name primary@127.0.0.1 --cookie ck -S mix run harness/reactive.exs
# -> results_reactive.csv

# RQ7 (kind): PDB sadar-kuorum menolak eviction pemecah-kuorum
bash code/k8s/rq7_eviction.sh        # -> data/results_eviction.csv

# Render semua figur dari data/*.csv
~/venv/bin/python code/analysis/plot.py     # -> figures/eval_*.pdf

# Tes
MIX_ENV=test mix test                                          # unit (policy)
MIX_ENV=test elixir --name primary@127.0.0.1 --cookie ck -S mix test --only cluster  # cluster
\end{lstlisting}
Atau pakai \code{Makefile} di akar proyek: \code{make test}, \code{make experiments},
\code{make figures}, \code{make rq7}, \code{make reproduce}.

% =============================================================================
\section{Tabel Pemetaan Ringkas (RQ $\rightarrow$ Skrip $\rightarrow$ CSV $\rightarrow$ Figur)}
\label{sec:peta}
\begin{longtable}{@{}p{0.06\linewidth} p{0.30\linewidth} p{0.30\linewidth} p{0.25\linewidth}@{}}
\toprule
\textbf{RQ} & \textbf{Skrip} & \textbf{CSV} & \textbf{Figur} \\
\midrule \endhead
RQ1 & \code{harness/run.exs} & \code{results\_safety.csv} & \code{eval\_safety.pdf} \\
RQ2 & \code{harness/run.exs} & \code{results\_efficiency.csv} & \code{eval\_efficiency.pdf} \\
RQ3 & \code{harness/policy.exs} & \code{results\_scale.csv} & \code{eval\_adaptivity.pdf} \\
RQ4 & \code{harness/policy.exs} & \code{results\_overhead.csv} & --- \\
RQ5 & \code{harness/policy.exs} & \code{results\_sensitivity.csv} & \code{eval\_sensitivity.pdf} \\
RQ6 & \code{harness/policy.exs} & \code{results\_scale.csv} & \code{eval\_scale.pdf} \\
RQ7 & \code{k8s/rq7\_eviction.sh} & \code{results\_eviction.csv} & --- \\
RQ8 & \code{harness/workload.exs} & \code{results\_workload.csv} & \code{eval\_workload.pdf} \\
RQ9 & \code{harness/reactive.exs} & \code{results\_reactive.csv} & \code{eval\_reactive.pdf} \\
RQ10 & \code{harness/policy.exs} & \code{results\_flap.csv} & \code{eval\_flap.pdf} \\
\bottomrule
\end{longtable}

% =============================================================================
\section{Glosarium Singkat}
\label{sec:glosarium}
\begin{description}[leftmargin=1.4em,style=nextline]
  \item[Kuorum (quorum)] Jumlah anggota hidup minimum (biasanya mayoritas) agar cluster tetap berwenang.
  \item[Split-brain] Cluster pecah jadi dua bagian yang sama-sama merasa berwenang $\rightarrow$ state
    bisa hilang saat pulih.
  \item[PodDisruptionBudget (PDB)] Aturan Kubernetes: paling sedikit \code{minAvailable} (atau paling
    banyak \code{maxUnavailable}) pod yang boleh down saat gangguan sukarela.
  \item[Gangguan sukarela] Pemeliharaan terencana: rolling update, \code{kubectl drain}, scale-down.
  \item[Eviction] Permintaan terkontrol untuk mengeluarkan sebuah pod; menghormati PDB (bisa ditolak).
  \item[Rolling update] Mengganti pod versi lama dengan versi baru secara bertahap.
  \item[GenServer] ``Proses server'' OTP yang memegang state dan menjawab pesan satu per satu.
  \item[Operator] Program yang terus mengamati cluster lalu menyesuaikan sesuatu (di sini: menambal PDB).
  \item[Ukuran diinginkan vs live] \emph{Diinginkan} = jumlah replika yang ditetapkan; \emph{live} =
    yang sedang terlihat. Kuorum dipatok ke yang \emph{diinginkan}.
  \item[libcluster / :peer / :rpc] Pembentuk cluster BEAM / node uji multi-proses / pemanggilan fungsi
    di node lain.
\end{description}

\end{document}
